% Nguồn: Tuyển tập đề tham khảo môn Toán TS10 2025-2026 – TP. Hồ Chí Minh
% Đợt 7: Quận 6 (Đề 2) và Quận 7 (Đề 1) (14 câu tự luận)

% ==============================================================================
% QUẬN 6 – ĐỀ 2
% ==============================================================================

\begin{bt}%[9D6H18-2]
% nguon: Đề tham khảo Tuyển sinh 10 TP.HCM 2025-2026 – Quận 6 – Đề 2 – Bài 1
Cho parabol $(P)\colon y = x^2$.
\begin{enumerate}[a)]
    \item Vẽ đồ thị $(P)$ trên mặt phẳng tọa độ $Oxy$.
    \item Tìm các điểm trên parabol (khác gốc tọa độ) cách đều hai trục tọa độ.
\end{enumerate}
\loigiai{
\begin{enumerate}[a)]
    \item Bảng giá trị tương ứng của $x$ và $y$:
    \begin{center}
    \begin{tabular}{|c|c|c|c|c|c|}
        \hline
        $x$ & $-2$ & $-1$ & $0$ & $1$ & $2$ \\ \hline
        $y = x^2$ & $4$ & $1$ & $0$ & $1$ & $4$ \\ \hline
    \end{tabular}
    \end{center}
    Đồ thị $(P)$ là một parabol đi qua gốc tọa độ $O(0;0)$, nhận trục tung $Oy$ làm trục đối xứng và bề lõm quay lên trên.
    \item Điểm $M(x; y) \in (P)$ cách đều hai trục tọa độ khi và chỉ khi:
    $|x| = |y|$.\\
    Vì $y = x^2 \ge 0$ nên ta có $|y| = y$, do đó:
    $y = |x| \Leftrightarrow x^2 = |x| \Leftrightarrow |x|^2 - |x| = 0 \Leftrightarrow |x|(|x| - 1) = 0 \Leftrightarrow \left[\begin{aligned} &|x| = 0 \\ &|x| = 1. \end{aligned}\right.$\\
    Vì điểm cần tìm khác gốc tọa độ $O(0;0)$ nên $|x| = 1 \Leftrightarrow x = \pm 1$.\\
    Với $x = 1 \Rightarrow y = 1$, ta được điểm $M_1(1; 1)$.\\
    Với $x = -1 \Rightarrow y = 1$, ta được điểm $M_2(-1; 1)$.\\
    Vậy có hai điểm trên parabol thỏa mãn yêu cầu bài toán là $(1; 1)$ và $(-1; 1)$.
\end{enumerate}
}
\end{bt}

\begin{bt}%[9D6H20-2]
% nguon: Đề tham khảo Tuyển sinh 10 TP.HCM 2025-2026 – Quận 6 – Đề 2 – Bài 2
Cho phương trình $5x^2 - 4x - 9 = 0$.
\begin{enumerate}[a)]
    \item Chứng minh phương trình trên có hai nghiệm phân biệt.
    \item Không giải phương trình, hãy tính giá trị của biểu thức: $B = x_1(x_1 - 3) + x_2(x_2 + 7x_1) - 3x_2$.
\end{enumerate}
\loigiai{
\begin{enumerate}[a)]
    \item Phương trình $5x^2 - 4x - 9 = 0$ có các hệ số $a = 5$ và $c = -9$.\\
    Vì $a \cdot c = 5 \cdot (-9) = -45 < 0$ nên phương trình luôn có hai nghiệm phân biệt $x_1, x_2$.
    \item Theo định lí Viète, ta có:
    $\left\{\begin{aligned} &S = x_1 + x_2 = \dfrac{4}{5} \\ &P = x_1 x_2 = -\dfrac{9}{5}. \end{aligned}\right.$\\
    Biến đổi biểu thức $B$:
    $\begin{aligned}
    B &= x_1(x_1 - 3) + x_2(x_2 + 7x_1) - 3x_2 \\
    &= x_1^2 - 3x_1 + x_2^2 + 7x_1 x_2 - 3x_2 \\
    &= (x_1^2 + x_2^2) + 7x_1 x_2 - 3(x_1 + x_2) \\
    &= \left[(x_1 + x_2)^2 - 2x_1 x_2\right] + 7x_1 x_2 - 3(x_1 + x_2) \\
    &= S^2 + 5P - 3S.
    \end{aligned}$\\
    Thay $S = \dfrac{4}{5}$ và $P = -\dfrac{9}{5}$ vào biểu thức, ta được:
    $B = \left(\dfrac{4}{5}\right)^2 + 5 \cdot \left(-\dfrac{9}{5}\right) - 3 \cdot \dfrac{4}{5} = \dfrac{16}{25} - 9 - \dfrac{12}{5} = \dfrac{16 - 225 - 60}{25} = -\dfrac{269}{25} = -10{,}76$.\\
    Vậy $B = -\dfrac{269}{25} = -10{,}76$.
\end{enumerate}
}
\end{bt}

\begin{bt}%[9D1H3-4]
% nguon: Đề tham khảo Tuyển sinh 10 TP.HCM 2025-2026 – Quận 6 – Đề 2 – Bài 3
Trên cánh đồng có diện tích $160\text{ ha}$ của một đơn vị sản xuất, người ta dành $60\text{ ha}$ để cấy thí điểm giống lúa mới, còn lại vẫn cấy giống lúa cũ. Khi thu hoạch, đầu tiên người ta gặt $8\text{ ha}$ giống lúa cũ và $7\text{ ha}$ giống lúa mới để đối chứng. Kết quả là $7\text{ ha}$ giống lúa mới cho thu hoạch nhiều hơn $8\text{ ha}$ giống lúa cũ là 2 tấn thóc. Biết rằng tổng số thóc (cả hai giống) thu hoạch cả vụ trên $160\text{ ha}$ là 860 tấn. Hỏi năng suất của mỗi giống lúa trên $1\text{ ha}$ là bao nhiêu tấn thóc?
\loigiai{
Diện tích cấy giống lúa cũ là:
$160 - 60 = 100\text{ (ha)}$.\\
Gọi $x$ (tấn/ha) là năng suất của giống lúa cũ và $y$ (tấn/ha) là năng suất của giống lúa mới ($x, y > 0$).\\
Theo đề bài, thu hoạch $7\text{ ha}$ giống lúa mới nhiều hơn $8\text{ ha}$ giống lúa cũ là 2 tấn thóc, ta có phương trình:
$7y - 8x = 2 \Leftrightarrow -8x + 7y = 2$.\\
Tổng sản lượng thóc thu hoạch trên cả cánh đồng $160\text{ ha}$ là 860 tấn, ta có phương trình:
$100x + 60y = 860 \Leftrightarrow 5x + 3y = 43$.\\
Ta có hệ phương trình:
$\left\{\begin{aligned} &-8x + 7y = 2 \\ &5x + 3y = 43 \end{aligned}\right. \Leftrightarrow \left\{\begin{aligned} &-24x + 21y = 6 \\ &35x + 21y = 301 \end{aligned}\right. \Leftrightarrow \left\{\begin{aligned} &59x = 295 \\ &3y = 43 - 5x \end{aligned}\right. \Leftrightarrow \left\{\begin{aligned} &x = 5 \\ &y = 6. \end{aligned}\right.$\\
Các giá trị $x = 5$ và $y = 6$ đều thỏa mãn điều kiện bài toán.\\
Vậy năng suất của giống lúa cũ là 5 tấn/ha và giống lúa mới là 6 tấn/ha.
}
\end{bt}

\begin{bt}%[9D6V21-3]
% nguon: Đề tham khảo Tuyển sinh 10 TP.HCM 2025-2026 – Quận 6 – Đề 2 – Bài 4
Một mảnh vườn hình chữ nhật có chiều rộng bằng $a\text{ (m)}$, chiều dài hơn chiều rộng $6\text{ (m)}$. Người ta làm lối đi xung quanh rộng $x\text{ (m)}$ (với $0 < 2x < a$).
\begin{enumerate}[a)]
    \item Viết biểu thức biểu thị diện tích phần còn lại của mảnh vườn theo $a$ và $x$.
    \item Tính diện tích phần còn lại khi $a = 30\text{ m}$, $x = 1\text{ m}$.
\end{enumerate}
\loigiai{
\begin{enumerate}[a)]
    \item Chiều rộng ban đầu của mảnh vườn là $a\text{ (m)}$.\\
    Chiều dài ban đầu của mảnh vườn là:
    $a + 6\text{ (m)}$.\\
    Vì lối đi xung quanh vườn có chiều rộng là $x\text{ (m)}$ nên:
    \begin{itemize}
        \item Chiều rộng phần đất còn lại là: $a - 2x\text{ (m)}$.
        \item Chiều dài phần đất còn lại là: $(a + 6) - 2x = a - 2x + 6\text{ (m)}$.
    \end{itemize}
    Biểu thức biểu thị diện tích phần đất còn lại của mảnh vườn là:
    $S = (a - 2x)(a - 2x + 6)\text{ (m}^2\text{)}$.
    \item Khi $a = 30\text{ m}$ và $x = 1\text{ m}$:\\
    Chiều rộng phần còn lại là:
    $30 - 2 \cdot 1 = 28\text{ (m)}$.\\
    Chiều dài phần còn lại là:
    $30 + 6 - 2 \cdot 1 = 34\text{ (m)}$.\\
    Diện tích phần còn lại của mảnh vườn là:
    $S = 28 \cdot 34 = 952\text{ (m}^2\text{)}$.
\end{enumerate}
}
\end{bt}

\begin{bt}%[9H0V31-3]
% nguon: Đề tham khảo Tuyển sinh 10 TP.HCM 2025-2026 – Quận 6 – Đề 2 – Bài 5
Một xe lu san đường (loại một trống lu) có đường kính trống lu là $0{,}96\text{ m}$ và chiều dài trống lu là $169\text{ cm}$. Người ta sử dụng loại xe lu này để làm phẳng một sân bóng đá hình chữ nhật có kích thước $120\text{ m} \times 90\text{ m}$. Cho rằng sân bóng cần được lăn 5 lần thì đạt tiêu chuẩn và mỗi trống lu chỉ lăn được tối đa với công suất $10\,000\text{ vòng/tuần}$. (Lấy $\pi \approx 3{,}14$).
\begin{enumerate}[a)]
    \item Tính tổng diện tích của mỗi xe lu có thể lăn trong 1 tuần (làm tròn đến hàng đơn vị mét vuông).
    \item Cần sử dụng ít nhất bao nhiêu xe lu để có thể hoàn thành công việc trong một tuần (biết rằng mỗi xe đều lăn hết công suất cho phép và các xe lu chỉ lăn trên phần sân riêng biệt)?
\end{enumerate}
\loigiai{
\begin{enumerate}[a)]
    \item Đổi chiều dài trống lu ra mét: $169\text{ cm} = 1{,}69\text{ m}$.\\
    Trống lu có dạng hình trụ với đường kính đáy $d = 0{,}96\text{ m}$ và chiều cao $h = 1{,}69\text{ m}$.\\
    Khi trống lu quay được 1 vòng, diện tích mặt sân được san phẳng đúng bằng diện tích xung quanh của hình trụ:
    $S_1 = \pi \cdot d \cdot h \approx 3{,}14 \cdot 0{,}96 \cdot 1{,}69 = 5{,}094528\text{ (m}^2\text{)}$.\\
    Trong 1 tuần, mỗi trống lu lăn tối đa $10\,000$ vòng nên tổng diện tích mỗi xe lu có thể lăn là:
    $S_{\text{tuần}} = 10\,000 \cdot S_1 \approx 10\,000 \cdot 5{,}094528 = 50\,945{,}28\text{ (m}^2\text{)} \approx 50\,945\text{ (m}^2\text{)}$.
    \item Diện tích sân bóng đá hình chữ nhật là:
    $S_{\text{sân}} = 120 \cdot 90 = 10\,800\text{ (m}^2\text{)}$.\\
    Vì sân bóng cần được lăn 5 lần nên tổng diện tích cần lăn là:
    $S_{\text{cần lăn}} = 10\,800 \cdot 5 = 54\,000\text{ (m}^2\text{)}$.\\
    Số lượng xe lu tối thiểu cần sử dụng là:
    $N = \left\lceil \dfrac{S_{\text{cần lăn}}}{S_{\text{tuần}}} \right\rceil = \left\lceil \dfrac{54\,000}{50\,945} \right\rceil = \lceil 1{,}06 \rceil = 2\text{ (xe lu)}$.\\
    Vậy cần sử dụng ít nhất 2 xe lu để hoàn thành công việc trong một tuần.
\end{enumerate}
}
\end{bt}

\begin{bt}%[9D8H26-3]
% nguon: Đề tham khảo Tuyển sinh 10 TP.HCM 2025-2026 – Quận 6 – Đề 2 – Bài 6
Số lượng máy tính để bàn và máy tính xách tay bán ra trong 6 tháng đầu năm của công ty X được thống kê như sau:
\begin{itemize}
    \item Tháng 1: 59 máy để bàn, 80 máy xách tay (chênh lệch 21 cái).
    \item Tháng 2: 43 máy để bàn, 70 máy xách tay (chênh lệch 27 cái).
    \item Tháng 3: 35 máy để bàn, 50 máy xách tay (chênh lệch 15 cái).
    \item Tháng 4: 24 máy để bàn, 60 máy xách tay (chênh lệch 36 cái).
    \item Tháng 5: 20 máy để bàn, 40 máy xách tay (chênh lệch 20 cái).
    \item Tháng 6: 14 máy để bàn, 50 máy xách tay (chênh lệch 36 cái).
\end{itemize}
\begin{enumerate}[a)]
    \item Trong 6 tháng đầu năm, tháng nào có sự chênh lệch giữa số lượng máy tính xách tay và máy tính để bàn được bán ra là ít nhất?
    \item Chọn ngẫu nhiên 1 tháng trong 6 tháng đầu năm, tính xác suất của các biến cố sau:
    \begin{itemize}
        \item $A$: “Tháng được chọn có số lượng máy tính để bàn mà công ty bán được không quá 30 cái”.
        \item $B$: “Tháng được chọn có số lượng máy tính xách tay mà công ty bán được ít nhất 50 cái”.
    \end{itemize}
\end{enumerate}
\loigiai{
\begin{enumerate}[a)]
    \item Chênh lệch giữa số lượng máy tính xách tay và máy tính để bàn qua các tháng:
    \begin{itemize}
        \item Tháng 1: $|80 - 59| = 21$ cái.
        \item Tháng 2: $|70 - 43| = 27$ cái.
        \item Tháng 3: $|50 - 35| = 15$ cái.
        \item Tháng 4: $|60 - 24| = 36$ cái.
        \item Tháng 5: $|40 - 20| = 20$ cái.
        \item Tháng 6: $|50 - 14| = 36$ cái.
    \end{itemize}
    Vậy Tháng 3 có sự chênh lệch ít nhất giữa số lượng máy tính xách tay và máy tính để bàn bán ra (chênh lệch 15 cái).
    \item Số phần tử của không gian mẫu khi chọn 1 trong 6 tháng là $n(\Omega) = 6$.
    \begin{itemize}
        \item Biến cố $A$: “Tháng được chọn có số lượng máy tính để bàn bán được không quá 30 cái” (tức $\le 30$ cái).\\
        Các tháng thỏa mãn là: Tháng 4 (24 cái), Tháng 5 (20 cái), Tháng 6 (14 cái).\\
        Số kết quả thuận lợi cho $A$ là $n(A) = 3$.\\
        Xác suất của biến cố $A$ là $P(A) = \dfrac{3}{6} = \dfrac{1}{2}$.
        \item Biến cố $B$: “Tháng được chọn có số lượng máy tính xách tay bán được ít nhất 50 cái” (tức $\ge 50$ cái).\\
        Các tháng thỏa mãn gồm: Tháng 1 (80 cái), Tháng 2 (70 cái), Tháng 3 (50 cái), Tháng 4 (60 cái), Tháng 6 (50 cái).\\
        Số kết quả thuận lợi cho $B$ là $n(B) = 5$.\\
        Xác suất của biến cố $B$ là $P(B) = \dfrac{5}{6}$.
    \end{itemize}
\end{enumerate}
}
\end{bt}

\begin{bt}%[9H9C29-3]
% nguon: Đề tham khảo Tuyển sinh 10 TP.HCM 2025-2026 – Quận 6 – Đề 2 – Bài 7
Cho đường tròn $(O, R)$ có đường kính $AB$. Gọi $H$ là điểm nằm giữa $O$ và $B$. Kẻ dây $CD$ vuông góc với $AB$ tại $H$. Trên cung nhỏ $AC$ lấy điểm $E$, kẻ $CK$ vuông góc với $AE$ tại $K$. Đường thẳng $DE$ cắt $CK$ tại $F$.
\begin{enumerate}[a)]
    \item Chứng minh tứ giác $AHCK$ là tứ giác nội tiếp.
    \item Chứng minh rằng $AH \cdot AB = AD^2$.
    \item Giả sử $H$ là trung điểm của $OB$. Tính độ dài cạnh $AF$ theo $R$.
\end{enumerate}
\loigiai{
\begin{enumerate}[a)]
    \item Vì $CD \perp AB$ tại $H$ nên $\widehat{AHC} = 90^\circ$.\\
    Vì $CK \perp AE$ tại $K$ nên $\widehat{AKC} = 90^\circ$.\\
    Xét tứ giác $AHCK$ có hai đỉnh kề $H$ và $K$ cùng nhìn cạnh $AC$ dưới một góc vuông, hay $\widehat{AHC} = \widehat{AKC} = 90^\circ$.\\
    Vậy tứ giác $AHCK$ nội tiếp đường tròn đường kính $AC$.
    \item Xét đường tròn $(O)$ đường kính $AB$, có điểm $D \in (O)$ nên $\widehat{ADB} = 90^\circ$ (góc nội tiếp chắn nửa đường tròn), do đó tam giác $ADB$ vuông tại $D$.\\
    Lại có $CD \perp AB$ tại $H$, suy ra $DH$ là đường cao của tam giác vuông $ADB$.\\
    Áp dụng hệ thức lượng trong tam giác vuông $ADB$, ta có:
    $AH \cdot AB = AD^2$.
    \item Xét tam giác $ACD$ cân tại $A$ vì đường kính $AB \perp CD$ tại $H$ nên $AB$ là đường trung trực của đoạn thẳng $CD$, suy ra $AC = AD$.\\
    Trong đường tròn $(O)$, tứ giác $ACDE$ nội tiếp nên $\widehat{FED} = \widehat{ACD}$.\\
    Lại có $CK \perp AE$ tại $K$ và $CD \perp AB$, suy ra tam giác $ACF$ cân tại $A$, do đó $AF = AC = AD$.\\
    Khi $H$ là trung điểm của $OB$:\\
    Ta có $OB = R \Rightarrow HB = \dfrac{R}{2}$.\\
    Độ dài đoạn thẳng $AH$ là:
    $AH = AB - HB = 2R - \dfrac{R}{2} = \dfrac{3R}{2}$.\\
    Áp dụng hệ thức ở câu b:
    $AD^2 = AH \cdot AB = \dfrac{3R}{2} \cdot 2R = 3R^2 \Rightarrow AD = R\sqrt{3}$.\\
    Vì $AF = AD$ nên ta có:
    $AF = R\sqrt{3}$.
\end{enumerate}
}
\end{bt}

% ==============================================================================
% QUẬN 7 – ĐỀ 1
% ==============================================================================

\begin{bt}%[9D6H18-2]
% nguon: Đề tham khảo Tuyển sinh 10 TP.HCM 2025-2026 – Quận 7 – Đề 1 – Bài 1
Cho parabol $(P)\colon y = \dfrac{x^2}{4}$.
\begin{enumerate}[a)]
    \item Vẽ đồ thị $(P)$ trên mặt phẳng tọa độ $Oxy$.
    \item Tìm những điểm $M$ thuộc $(P)$ có hoành độ và tung độ là hai số đối nhau.
\end{enumerate}
\loigiai{
\begin{enumerate}[a)]
    \item Bảng giá trị tương ứng của $x$ và $y$:
    \begin{center}
    \begin{tabular}{|c|c|c|c|c|c|}
        \hline
        $x$ & $-4$ & $-2$ & $0$ & $2$ & $4$ \\ \hline
        $y = \dfrac{x^2}{4}$ & $4$ & $1$ & $0$ & $1$ & $4$ \\ \hline
    \end{tabular}
    \end{center}
    Đồ thị $(P)$ là một parabol đi qua gốc tọa độ $O(0;0)$, nhận trục tung $Oy$ làm trục đối xứng và bề lõm quay lên trên.
    \item Điểm $M(x_M; y_M) \in (P)$ có hoành độ và tung độ đối nhau, tức là $y_M = -x_M$.\\
    Vì $M \in (P)$ nên ta có phương trình:
    $\dfrac{x_M^2}{4} = -x_M \Leftrightarrow x_M^2 + 4x_M = 0 \Leftrightarrow x_M(x_M + 4) = 0 \Leftrightarrow \left[\begin{aligned} &x_M = 0 \\ &x_M = -4. \end{aligned}\right.$\\
    Với $x_M = 0 \Rightarrow y_M = 0$, ta được điểm $O(0; 0)$ (số 0 là số đối của chính nó).\\
    Với $x_M = -4 \Rightarrow y_M = 4$, ta được điểm $M(-4; 4)$.\\
    Vậy các điểm cần tìm là $O(0; 0)$ và $M(-4; 4)$.
\end{enumerate}
}
\end{bt}

\begin{bt}%[9D6H20-2]
% nguon: Đề tham khảo Tuyển sinh 10 TP.HCM 2025-2026 – Quận 7 – Đề 1 – Bài 2
Cho phương trình $2x^2 - 6x - 1 = 0$ có hai nghiệm phân biệt là $x_1, x_2$.
\begin{enumerate}[a)]
    \item Chứng minh phương trình trên có hai nghiệm phân biệt.
    \item Không giải phương trình, hãy tính giá trị của biểu thức: $A = \dfrac{x_1 - 2}{x_2 - 1} + \dfrac{x_2 - 2}{x_1 - 1}$.
\end{enumerate}
\loigiai{
\begin{enumerate}[a)]
    \item Phương trình $2x^2 - 6x - 1 = 0$ có $a = 2$ và $c = -1$.\\
    Vì $a \cdot c = 2 \cdot (-1) = -2 < 0$ nên phương trình luôn có hai nghiệm phân biệt $x_1, x_2$.
    \item Theo định lí Viète, ta có:
    $\left\{\begin{aligned} &S = x_1 + x_2 = \dfrac{6}{2} = 3 \\ &P = x_1 x_2 = -\dfrac{1}{2}. \end{aligned}\right.$\\
    Quy đồng mẫu thức biểu thức $A$:
    $\begin{aligned}
    A &= \dfrac{(x_1 - 2)(x_1 - 1) + (x_2 - 2)(x_2 - 1)}{(x_2 - 1)(x_1 - 1)} \\
    &= \dfrac{(x_1^2 - 3x_1 + 2) + (x_2^2 - 3x_2 + 2)}{x_1 x_2 - (x_1 + x_2) + 1} \\
    &= \dfrac{(x_1^2 + x_2^2) - 3(x_1 + x_2) + 4}{x_1 x_2 - (x_1 + x_2) + 1} \\
    &= \dfrac{\left[(x_1 + x_2)^2 - 2x_1 x_2\right] - 3(x_1 + x_2) + 4}{x_1 x_2 - (x_1 + x_2) + 1} \\
    &= \dfrac{S^2 - 2P - 3S + 4}{P - S + 1}.
    \end{aligned}$\\
    Thay $S = 3$ và $P = -\dfrac{1}{2}$ vào biểu thức:
    $\begin{aligned}
    \text{Tử số} &= 3^2 - 2 \cdot \left(-\dfrac{1}{2}\right) - 3 \cdot 3 + 4 = 9 + 1 - 9 + 4 = 5, \\
    \text{Mẫu số} &= -\dfrac{1}{2} - 3 + 1 = -\dfrac{5}{2}.
    \end{aligned}$\\
    Do đó:
    $A = \dfrac{5}{-\dfrac{5}{2}} = -2$.\\
    Vậy $A = -2$.
\end{enumerate}
}
\end{bt}

\begin{bt}%[9D8H26-3]
% nguon: Đề tham khảo Tuyển sinh 10 TP.HCM 2025-2026 – Quận 7 – Đề 1 – Bài 3
Tấm bìa cứng A hình tròn được chia thành 3 hình quạt bằng nhau, đánh số 1; 2; 3 và tấm bìa cứng B hình tròn được chia thành 5 hình quạt bằng nhau, đánh số 1; 2; 3; 4; 5. Trục quay của A và B được gắn mũi tên ở tâm. Bạn Bình quay tấm bìa A, bạn An quay tấm bìa B. Quan sát xem mũi tên dừng ở hình quạt nào trên hai tấm bìa.
\begin{enumerate}[a)]
    \item Mô tả không gian mẫu của phép thử.
    \item Tính xác suất của các biến cố sau:
    \begin{itemize}
        \item $T$: “Tích hai số ở hình quạt mà hai mũi tên chỉ vào bằng 6”.
        \item $M$: “Tích hai số ở hình quạt mà hai mũi tên chỉ vào nhỏ hơn 5”.
        \item $L$: “Tích hai số ở hình quạt mà hai mũi tên chỉ vào là số chẵn”.
    \end{itemize}
\end{enumerate}
\loigiai{
\begin{enumerate}[a)]
    \item Kí hiệu kết quả phép thử là cặp số $(a; b)$, trong đó $a \in \{1; 2; 3\}$ là số trên tấm bìa A và $b \in \{1; 2; 3; 4; 5\}$ là số trên tấm bìa B.\\
    Không gian mẫu của phép thử là tập hợp gồm 15 phần tử:
    $\Omega = \{(a; b) \mid a \in \{1; 2; 3\}, b \in \{1; 2; 3; 4; 5\}\}$.\\
    Số phần tử của không gian mẫu là $n(\Omega) = 3 \cdot 5 = 15$.
    \item Tính xác suất của các biến cố:
    \begin{itemize}
        \item Biến cố $T$: “Tích hai số bằng 6” ($a \cdot b = 6$).\\
        Các cặp số có tích bằng 6 là: $(2; 3)$ và $(3; 2)$.\\
        Số kết quả thuận lợi là $n(T) = 2$.\\
        Xác suất của biến cố $T$ là:
        $P(T) = \dfrac{2}{15}$.
        \item Biến cố $M$: “Tích hai số nhỏ hơn 5” ($a \cdot b < 5$).\\
        Các kết quả thỏa mãn gồm:
        \begin{itemize}
            \item Với $a = 1$: $b \in \{1; 2; 3; 4\}$ (có 4 cặp: $(1; 1), (1; 2), (1; 3), (1; 4)$).
            \item Với $a = 2$: $b \in \{1; 2\}$ (có 2 cặp: $(2; 1), (2; 2)$).
            \item Với $a = 3$: $b \in \{1\}$ (có 1 cặp: $(3; 1)$).
        \end{itemize}
        Số kết quả thuận lợi là $n(M) = 4 + 2 + 1 = 7$.\\
        Xác suất của biến cố $M$ là:
        $P(M) = \dfrac{7}{15}$.
        \item Biến cố $L$: “Tích hai số là số chẵn”.\\
        Biến cố đối $\overline{L}$ là “Tích hai số là số lẻ”, xảy ra khi cả hai số $a$ và $b$ đều là số lẻ.\\
        Số cách chọn $a$ lẻ từ $\{1; 3\}$ là 2 cách.\\
        Số cách chọn $b$ lẻ từ $\{1; 3; 5\}$ là 3 cách.\\
        Số kết quả để tích là số lẻ là $n(\overline{L}) = 2 \cdot 3 = 6$.\\
        Số kết quả thuận lợi cho biến cố $L$ là:
        $n(L) = 15 - 6 = 9$.\\
        Xác suất của biến cố $L$ là:
        $P(L) = \dfrac{9}{15} = \dfrac{3}{5}$.
    \end{itemize}
\end{enumerate}
}
\end{bt}

\begin{bt}%[9D6V21-3]
% nguon: Đề tham khảo Tuyển sinh 10 TP.HCM 2025-2026 – Quận 7 – Đề 1 – Bài 4
Người ta làm một lối đi theo chiều dài và chiều rộng của một sân cỏ hình chữ nhật. Em hãy tính chiều rộng $x$ của lối đi, biết rằng lối đi có diện tích bằng $46\text{ m}^2$, sân cỏ ban đầu có chiều dài $15\text{ m}$ và chiều rộng $6\text{ m}$ (với $0 < x < 6$).
\begin{enumerate}[a)]
    \item Viết biểu thức tính diện tích của lối đi theo $x$.
    \item Tìm $x$.
\end{enumerate}
\loigiai{
\begin{enumerate}[a)]
    \item Diện tích toàn bộ sân cỏ ban đầu là:
    $S_0 = 15 \cdot 6 = 90\text{ (m}^2\text{)}$.\\
    Khi bớt lối đi rộng $x\text{ (m)}$ dọc theo một chiều dài và một chiều rộng của sân:
    \begin{itemize}
        \item Chiều dài phần sân cỏ còn lại là: $15 - x\text{ (m)}$.
        \item Chiều rộng phần sân cỏ còn lại là: $6 - x\text{ (m)}$.
    \end{itemize}
    Diện tích phần sân cỏ còn lại là:
    $S_{\text{còn}} = (15 - x)(6 - x) = x^2 - 21x + 90\text{ (m}^2\text{)}$.\\
    Diện tích của lối đi là hiệu diện tích sân ban đầu và phần còn lại:
    $S_{\text{lối đi}} = S_0 - S_{\text{còn}} = 90 - (x^2 - 21x + 90) = 21x - x^2\text{ (m}^2\text{)}$.
    \item Theo đề bài, diện tích lối đi bằng $46\text{ m}^2$, ta có phương trình:
    $21x - x^2 = 46 \Leftrightarrow x^2 - 21x + 46 = 0 \Leftrightarrow (x - 2)(x - 23) = 0 \Leftrightarrow \left[\begin{aligned} &x = 2\text{ (thỏa mãn } 0 < x < 6) \\ &x = 23\text{ (loại vì } x > 6). \end{aligned}\right.$\\
    Vậy chiều rộng của lối đi là $x = 2\text{ m}$.
\end{enumerate}
}
\end{bt}

\begin{bt}%[9H0V31-3]
% nguon: Đề tham khảo Tuyển sinh 10 TP.HCM 2025-2026 – Quận 7 – Đề 1 – Bài 5
Một cái ly thủy tinh, phần phía trên là hình nón có chiều cao $7\text{ cm}$, đáy là đường tròn bán kính $4\text{ cm}$. Thể tích hình nón được tính theo công thức $V = \dfrac{1}{3}\pi r^2 h$.
\begin{enumerate}[a)]
    \item Tính thể tích của cái ly (bề dày của ly không đáng kể, lấy $\pi \approx 3{,}14$).
    \item Biết trong ly đang chứa rượu với mức rượu cách miệng ly là $3\text{ cm}$. Hỏi thể tích còn lại của ly rượu chiếm bao nhiêu phần của thể tích ly (làm tròn đến hai chữ số thập phân)?
\end{enumerate}
\loigiai{
\begin{enumerate}[a)]
    \item Thể tích của cái ly hình nón là:
    $V = \dfrac{1}{3}\pi r^2 h \approx \dfrac{1}{3} \cdot 3{,}14 \cdot 4^2 \cdot 7 = \dfrac{351{,}68}{3} \approx 117{,}23\text{ (cm}^3\text{)}$.
    \item Mức rượu cách miệng ly là $3\text{ cm}$, suy ra chiều cao của phần hình nón chứa rượu là:
    $h' = 7 - 3 = 4\text{ (cm)}$.\\
    Phần hình nón chứa rượu đồng dạng với phần hình nón của toàn bộ cái ly với tỉ số đồng dạng:
    $k = \dfrac{h'}{h} = \dfrac{4}{7}$.\\
    Do đó, tỉ số giữa thể tích rượu và thể tích toàn bộ cái ly là:
    $\dfrac{V_{\text{rượu}}}{V} = k^3 = \left(\dfrac{4}{7}\right)^3 = \dfrac{64}{343}$.\\
    Thể tích còn lại (phần không chứa rượu) chiếm số phần thể tích ly là:
    $1 - \dfrac{V_{\text{rượu}}}{V} = 1 - \dfrac{64}{343} = \dfrac{279}{343} \approx 0{,}81$.
\end{enumerate}
}
\end{bt}

\begin{bt}%[9D1H3-4]
% nguon: Đề tham khảo Tuyển sinh 10 TP.HCM 2025-2026 – Quận 7 – Đề 1 – Bài 6
Thép không gỉ Ferritic là họ thép hợp kim có chứa từ $12\%$ đến $27\%$ crôm. Một nhà máy luyện thép hiện có sẵn một lượng hợp kim thép chứa $10\%$ crôm và một lượng hợp kim thép chứa $30\%$ crôm (các nguyên liệu không bị hao hụt trong quá trình luyện).
\begin{enumerate}[a)]
    \item Tính khối lượng hợp kim thép mỗi loại từ hai loại thép trên dùng để luyện được $500$ tấn thép chứa $16\%$ crôm.
    \item Nhà máy dự định luyện ra loại thép không gỉ Ferritic từ $100$ tấn thép chứa $10\%$ crôm và $x$ tấn thép chứa $30\%$ crôm. Hỏi $x$ nằm trong khoảng nào?
\end{enumerate}
\loigiai{
\begin{enumerate}[a)]
    \item Gọi $x, y$ lần lượt là khối lượng hợp kim thép chứa $10\%$ crôm và chứa $30\%$ crôm cần dùng ($x, y > 0$, đơn vị tấn).\\
    Tổng khối lượng thép cần luyện là $500$ tấn, ta có phương trình:
    $x + y = 500$.\\
    Khối lượng crôm nguyên chất trong $500$ tấn thép chứa $16\%$ crôm là:
    $500 \cdot 16\% = 80\text{ (tấn)}$.\\
    Theo bài ra ta có phương trình:
    $10\% x + 30\% y = 80 \Leftrightarrow 0{,}1x + 0{,}3y = 80 \Leftrightarrow x + 3y = 800$.\\
    Ta có hệ phương trình:
    $\left\{\begin{aligned} &x + y = 500 \\ &x + 3y = 800 \end{aligned}\right. \Leftrightarrow \left\{\begin{aligned} &2y = 300 \\ &x = 500 - y \end{aligned}\right. \Leftrightarrow \left\{\begin{aligned} &y = 150 \\ &x = 350. \end{aligned}\right.$\\
    Vậy cần dùng $350$ tấn thép chứa $10\%$ crôm và $150$ tấn thép chứa $30\%$ crôm.
    \item Khi trộn $100$ tấn thép chứa $10\%$ crôm với $x$ tấn thép chứa $30\%$ crôm ($x > 0$):\\
    Tổng khối lượng hợp kim thu được là: $100 + x$ (tấn).\\
    Tổng khối lượng crôm nguyên chất là:
    $100 \cdot 10\% + x \cdot 30\% = 10 + 0{,}3x$ (tấn).\\
    Tỉ lệ phần trăm crôm trong hợp kim mới là:
    $\dfrac{10 + 0{,}3x}{100 + x}$.\\
    Để hợp kim mới là thép không gỉ Ferritic (chứa từ $12\%$ đến $27\%$ crôm), ta có bất phương trình kép:
    $12\% \le \dfrac{10 + 0{,}3x}{100 + x} \le 27\% \Leftrightarrow 0{,}12 \le \dfrac{10 + 0{,}3x}{100 + x} \le 0{,}27$.\\
    Vì $100 + x > 0$ nên:
    $\left\{\begin{aligned} &0{,}12(100 + x) \le 10 + 0{,}3x \\ &10 + 0{,}3x \le 0{,}27(100 + x) \end{aligned}\right. \Leftrightarrow \left\{\begin{aligned} &12 + 0{,}12x \le 10 + 0{,}3x \\ &10 + 0{,}3x \le 27 + 0{,}27x \end{aligned}\right. \Leftrightarrow \left\{\begin{aligned} &0{,}18x \ge 2 \\ &0{,}03x \le 17 \end{aligned}\right. \Leftrightarrow \left\{\begin{aligned} &x \ge \dfrac{100}{9} \approx 11{,}11 \\ &x \le \dfrac{1700}{3} \approx 566{,}67. \end{aligned}\right.$\\
    Vậy giá trị của $x$ nằm trong khoảng từ $\dfrac{100}{9}$ tấn đến $\dfrac{1700}{3}$ tấn (xấp xỉ từ $11{,}11$ tấn đến $566{,}67$ tấn).
\end{enumerate}
}
\end{bt}

\begin{bt}%[9H9C29-3]
% nguon: Đề tham khảo Tuyển sinh 10 TP.HCM 2025-2026 – Quận 7 – Đề 1 – Bài 7
Cho tam giác $ABC$ có ba góc nhọn ($AB < AC$). Đường tròn tâm $O$ đường kính $BC$ cắt hai cạnh $AB, AC$ lần lượt tại $E$ và $F$ ($E \ne B, F \ne C$). Các đoạn thẳng $BF$ và $CE$ cắt nhau tại $H$, tia $AH$ cắt $BC$ tại $K$.
\begin{enumerate}[a)]
    \item Chứng minh $\widehat{BEC} = \widehat{BFC} = 90^\circ$, từ đó suy ra tứ giác $AEHF$ nội tiếp.
    \item Gọi $D$ là giao điểm của $AH$ và $(O)$ ($D$ nằm giữa $A$ và $H$). Chứng minh $BD^2 = BK \cdot BC$ và $\widehat{BDH} = \widehat{BFD}$.
    \item Trong trường hợp $\widehat{BAC} = 60^\circ$ và $BC = 6\text{ cm}$, tính độ dài đoạn thẳng $EF$ và bán kính đường tròn ngoại tiếp tam giác $AEF$.
\end{enumerate}
\loigiai{
\begin{enumerate}[a)]
    \item Xét đường tròn $(O)$ đường kính $BC$ có các điểm $E, F \in (O)$ nên:
    $\widehat{BEC} = 90^\circ$ và $\widehat{BFC} = 90^\circ$ (các góc nội tiếp chắn nửa đường tròn).\\
    Suy ra $CE \perp AB$ và $BF \perp AC$, do đó $CE$ và $BF$ là hai đường cao của tam giác $ABC$.\\
    Xét tứ giác $AEHF$ có tổng hai góc đối diện:
    $\widehat{AEH} + \widehat{AFH} = 90^\circ + 90^\circ = 180^\circ$.\\
    Vậy tứ giác $AEHF$ nội tiếp đường tròn đường kính $AH$.
    \item Vì $H$ là giao điểm của hai đường cao $BF$ và $CE$ trong tam giác $ABC$ nên $H$ là trực tâm của tam giác $ABC$, suy ra $AK \perp BC$ tại $K$.\\
    Vì $D \in (O)$ đường kính $BC$ nên tam giác $BDC$ vuông tại $D$, có $DK$ là đường cao ($AK \perp BC$ tại $K$).\\
    Áp dụng hệ thức lượng trong tam giác vuông $BDC$:
    $BD^2 = BK \cdot BC$.\\
    Xét hai tam giác vuông $BKH$ và $BFC$ có góc $\widehat{KBF}$ chung, suy ra $\triangle BKH \sim \triangle BFC$ (g-g).\\
    Do đó $\dfrac{BK}{BF} = \dfrac{BH}{BC} \Leftrightarrow BK \cdot BC = BH \cdot BF$.\\
    Kết hợp với hệ thức trên, ta có $BD^2 = BH \cdot BF \Leftrightarrow \dfrac{BD}{BF} = \dfrac{BH}{BD}$.\\
    Xét $\triangle BDH$ và $\triangle BFD$ có:
    $\left\{\begin{aligned} &\widehat{DBF}\text{ là góc chung} \\ &\dfrac{BD}{BF} = \dfrac{BH}{BD}\text{ (chứng minh trên)}. \end{aligned}\right.$\\
    Do đó $\triangle BDH \sim \triangle BFD$ (c-g-c), suy ra $\widehat{BDH} = \widehat{BFD}$.
    \item Khi $\widehat{BAC} = 60^\circ$ và $BC = 6\text{ cm}$:\\
    Xét tứ giác $BCEF$ có $\widehat{BEC} = \widehat{BFC} = 90^\circ$ nên $BCEF$ nội tiếp đường tròn $(O)$.\\
    Suy ra $\widehat{AFE} = \widehat{ABC}$, do đó $\triangle AEF \sim \triangle ACB$ (g-g).\\
    Tỉ số đồng dạng là:
    $k = \dfrac{AE}{AC} = \cos \widehat{BAC} = \cos 60^\circ = \dfrac{1}{2}$.\\
    Từ đó ta có:
    $EF = k \cdot BC = \cos 60^\circ \cdot BC = \dfrac{1}{2} \cdot 6 = 3\text{ (cm)}$.\\
    Tính bán kính đường tròn ngoại tiếp tam giác $AEF$:\\
    Theo câu a, tứ giác $AEHF$ nội tiếp đường tròn đường kính $AH$, nên đường tròn ngoại tiếp tam giác $AEF$ chính là đường tròn đường kính $AH$, có bán kính là $R' = \dfrac{AH}{2}$.\\
    Gọi $R_0$ là bán kính đường tròn ngoại tiếp tam giác $ABC$.\\
    Theo định lí sin trong tam giác $ABC$:
    $2R_0 = \dfrac{BC}{\sin \widehat{BAC}} = \dfrac{6}{\sin 60^\circ} = \dfrac{6}{\dfrac{\sqrt{3}}{2}} = 4\sqrt{3}\text{ (cm)} \Rightarrow R_0 = 2\sqrt{3}\text{ (cm)}$.\\
    Độ dài đoạn thẳng nối từ đỉnh $A$ đến trực tâm $H$ được tính theo công thức:
    $AH = 2R_0 \cos \widehat{BAC} = 2 \cdot (2\sqrt{3}) \cdot \cos 60^\circ = 4\sqrt{3} \cdot \dfrac{1}{2} = 2\sqrt{3}\text{ (cm)}$.\\
    Vậy bán kính đường tròn ngoại tiếp tam giác $AEF$ là:
    $R' = \dfrac{AH}{2} = \dfrac{2\sqrt{3}}{2} = \sqrt{3}\text{ (cm)} \approx 1{,}73\text{ (cm)}$.
\end{enumerate}
}
\end{bt}
